\documentclass[a4paper,12pt]{article}
\usepackage[utf8]{inputenc}
\usepackage[spanish]{babel}
\usepackage{amsmath}
\usepackage{amsfonts}
\usepackage{amssymb}
\usepackage{graphicx}
\usepackage{siunitx}
\usepackage{tikz}
\usetikzlibrary{babel}
\usepackage{geometry}
\usepackage{bookmark}
\usepackage{enumitem}

\geometry{margin=2.5cm}

\title{Hoja de Trabajo: Física I - Conceptos Fundamentales}
\author{Arquitecto: \underline{\hspace{5cm}} \quad Fecha: \underline{\hspace{3cm}}}
\date{}

\begin{document}

\maketitle

\section{Parte I: Teoría (10 puntos)}
Responda de forma clara y precisa a las siguientes interrogantes:

\begin{enumerate}[itemsep=1.5em]
    \item \textbf{Etimología}: ¿De qué vocablo griego proviene la palabra ``Física'' y qué significa?
    \item \textbf{Definición}: ¿Qué estudia la ciencia de la Física?
    \item \textbf{Divisiones}: Mencione las dos grandes divisiones de la física y dé tres ejemplos de ramas de la Física Clásica.
    \item \textbf{S.I.}: Enumere cinco (5) magnitudes fundamentales del Sistema Internacional junto con sus unidades básicas.
    \item \textbf{Derivadas}: Explique qué son las magnitudes derivadas y proporcione dos ejemplos incluyendo sus unidades.
    \item \textbf{Interdisciplinariedad}: Mencione tres ciencias con las que la física tiene una relación directa.
    \item \textbf{Cifras Significativas (Ceros)}: Explique la regla de los ceros situados a la izquierda y la de los ceros intermedios.
    \item \textbf{Redondeo (Suma/Resta)}: ¿Qué criterio define la cantidad de decimales en el resultado de una suma o resta?
    \item \textbf{Redondeo (Prod/Div)}: ¿Cuál es la regla para determinar las cifras significativas en una multiplicación o división?
    \item \textbf{Notación Científica}: Describa la principal utilidad de este sistema de representación.
\end{enumerate}

\newpage

\section{Parte II: Ejercicios Prácticos}

\subsection{Cifras Significativas y Redondeo}
1. Determine el número de cifras significativas en:
\begin{itemize}
    \item a) 124 \hfill \underline{\hspace{2cm}}
    \item b) 0.0035 \hfill \underline{\hspace{2cm}}
    \item c) 12.0 \hfill \underline{\hspace{2cm}}
    \item d) $1.2 \times 10^3$ \hfill \underline{\hspace{2cm}}
\end{itemize}

2. Realice las siguientes operaciones aplicando las reglas de redondeo:
\begin{itemize}
    \item Suma: $3.95 + 4.198 + 12.17 =$ \underline{\hspace{3cm}}
    \item Resta: $18.94 - 8.4 =$ \underline{\hspace{3cm}}
    \item Multiplicación: $18.94 \times 12.713 =$ \underline{\hspace{3cm}}
    \item División: $36.72 / 4.5 =$ \underline{\hspace{3cm}}
\end{itemize}

\subsection{Notación Científica}
3. Exprese en notación científica o decimal según corresponda:
\begin{itemize}
    \item 900,000 $\rightarrow$ \underline{\hspace{4cm}}
    \item 0.000000658 $\rightarrow$ \underline{\hspace{4cm}}
    \item $1.03 \times 10^{-5} \rightarrow$ \underline{\hspace{4cm}}
    \item $7.00 \times 10^{12} \rightarrow$ \underline{\hspace{4cm}}
\end{itemize}

4. Realice las operaciones:
\begin{itemize}
    \item $(3.2 \times 10^6) \times (2.19 \times 10^8) =$ \underline{\hspace{4cm}}
    \item $(2.13 \times 10^9) / (6.98 \times 10^{11}) =$ \underline{\hspace{4cm}}
\end{itemize}

\subsection{Conversiones Avanzadas}
Resuelva las siguientes conversiones (muestre el procedimiento):
\begin{enumerate}
    \item Convierta 6 km a pies.
    \item Convierta 5 millas/h a m/s.
    \item Convierta $96,500 \text{ cm}^3/\text{min}$ a gal/s.
    \item Convierta 1.2 km a pulgadas (in).
    \item Convierta 0.94 gal/s a $\text{cm}^3/\text{hr}$.
    \item Convierta 1500 km/s a millas/min.
\end{enumerate}

\newpage

\section{Parte III: Estática y Vectores}

\subsection{Ejercicio \#8: Dirección del vector resultante}
\begin{center}
\begin{tikzpicture}[scale=0.5]
    \draw[dashed, ->] (-8,0) -- (10,0) node[right] {$x$};
    \draw[dashed, ->] (0,-8) -- (0,6) node[above] {$y$};
    \draw[ultra thick, ->, blue] (0,0) -- ({10*cos(143)},{10*sin(143)}) node[left] {10u};
    \draw[ultra thick, ->, red] (0,0) -- ({4*cos(30)},{4*sin(30)}) node[right] {4u};
    \draw[ultra thick, ->, black] (0,0) -- (8,0) node[below] {8u};
    \draw[ultra thick, ->, gray] (0,0) -- (0,-6) node[right] {6u};
    \draw (0.8,0) arc (0:30:0.8) node[midway, right] {\tiny $30^\circ$};
    \draw (0,1.2) arc (90:143:1.2) node[midway, above left] {\tiny $53^\circ$};
\end{tikzpicture}
\end{center}
Opciones: \textbf{A) 30$^\circ$ \quad B) 37$^\circ$ \quad C) 45$^\circ$ \quad D) 53$^\circ$ \quad E) 60$^\circ$}

\subsection{Ejercicio \#9: Módulo del vector resultante}
\begin{center}
\begin{tikzpicture}[scale=0.5]
    \draw[->] (-5,0) -- (14,0) node[right] {$x$};
    \draw[->] (0,-8) -- (0,5) node[above] {$y$};
    \draw[ultra thick, ->, blue] (0,0) -- ({4*cos(135)},{4*sin(135)}) node[above left] {$2\sqrt{2}$};
    \draw[ultra thick, ->, red] (0,0) -- (13,0) node[below] {13};
    \draw[ultra thick, ->, orange] (0,0) -- ({10*cos(233)},{10*sin(233)}) node[below left] {10};
    \draw (-0.8,-0.1) arc (180:233:0.8) node[midway, left] {\tiny $53^\circ$};
    \draw (0.8,0.1) arc (0:135:0.8) node[midway, above right] {\tiny $45^\circ$};
\end{tikzpicture}
\end{center}
Opciones: \textbf{A) 1 \quad B) 2 \quad C) 3 \quad D) 4 \quad E) 5}

\subsection{Ejercicio \#14: Cálculo de Fuerza ``F''}
Si la fuerza resultante es horizontal, halle ``F''.
\begin{center}
\begin{tikzpicture}[scale=0.4]
    \draw[->] (-10,0) -- (10,0);
    \draw[->] (0,-10) -- (0,5);
    \draw[ultra thick, ->, orange] (0,0) -- ({8*cos(53)},{8*sin(53)}) node[right] {F};
    \draw[ultra thick, ->] (0,0) -- ({14*cos(210)},{14*sin(210)}) node[left] {14u};
    \draw[ultra thick, ->] (0,0) -- (0,-9) node[right] {9u};
    \draw (-1.5,0) arc (180:210:1.5) node[midway, left] {\tiny $30^\circ$};
    \draw (1.5,0) arc (0:53:1.5) node[midway, right] {\tiny $53^\circ$};
\end{tikzpicture}
\end{center}
Opciones: \textbf{A) 10N \quad B) 20 \quad C) 30 \quad D) 40 \quad E) 20$\sqrt{3}$}

\subsection{Ejercicio \#15: Módulo de la resultante}
\begin{center}
\begin{tikzpicture}[scale=0.4]
    \draw[<->] (-12,0) -- (12,0);
    \draw[<->] (0,-12) -- (0,12);
    \draw[ultra thick, ->] (0,0) -- ({10*cos(101)},{10*sin(101)}) node[above] {50u};
    \draw[ultra thick, ->] (0,0) -- ({10*cos(350)},{10*sin(350)}) node[right] {50u};
    \draw (0,1.5) arc (90:101:1.5) node[midway, above] {\tiny $11^\circ$};
    \draw (1.5,0) arc (0:-10:1.5) node[midway, right] {\tiny $10^\circ$};
\end{tikzpicture}
\end{center}
Opciones: \textbf{A) 8$\sqrt{10}$u \quad B) 18$\sqrt{10}$ \quad C) 50$\sqrt{2}$ \quad D) 50 \quad E) 40$\sqrt{5}$}

\end{document}
